\section{Related Work}

\subsection{Scaling Laws for Language Models} Scaling laws have proven to be a powerful framework for understanding and predicting the performance of language models. Existing studies, such as \citet{kaplan_scale} and \citet{hoffmann2022training}, reveal that power-law relationships govern model performance as a function of factors like model size, data size, and compute budget, offering predictable performance improvements with increased resources.

\citet{hoffmann2022training} emphasizes the critical balance between model size and the number of training tokens when the training compute budget is fixed, showing that scaling the model without corresponding data increases can lead to suboptimal performance. Additionally, \citet{Bi2024DeepSeekLS} explores more nuanced scaling behaviors by incorporating data quality, demonstrating that higher-quality data allows for more efficient scaling, and thus, a larger portion of the compute budget should be allocated to increasing model size.

Recent work extends scaling law analysis to specialized contexts, including over-training~\citep{Gadre2024overtrain}, downstream task performance, and multilingual or multi-modal settings, where scaling laws provide valuable insights and can be adapted to address specific challenges.

\subsection{Scaling Laws for MoEs} 
Mixture-of-Experts (MoE) models~\citep{shazeer2017,lepikhin2021gshard,fedus_switch,deepseekr1} have emerged as a powerful architecture for language modeling, primarily because they decouple computational cost from parameter count. This separation between parameters and \flops per token in \moe architectures calls for scaling laws that can accurately factor in the contributions of both.

Previous research on the scaling behavior of MoE models has established foundational scaling laws, incorporating factors such as total parameter count, the number of experts, and the granularity of these experts~\citep{clark2022unified, ludziejewski2024scaling, wang-etal-2024-scaling}. However, these studies typically assume a fixed configuration for other critical variables influencing \flops per token, such as the number of active experts per input. In contrast, we propose a generalized scaling law that considers variables like active parameter count and sparsity level, thereby expanding the applicability of MoE scaling laws.

A common theme in the literature suggests that training sparser models—achieved by increasing the number of smaller experts—offers significant gains in efficiency for both pretraining and inference phases. Through a comprehensive large-scale study, we provide empirical evidence for this, analyzing the impact of sparsity level on efficiency and defining optimal configurations.

Supporting this, \citet{Du2021GLaMES} demonstrates GLaM’s superior efficiency and performance compared to GPT-3, showing that MoE architectures can achieve high performance with significantly lower computational and energy costs. Further insights are offered by \citet{clark2022unified}, who analyze scaling behaviors across various MoE routing techniques. While their study finds that MoEs generally outperform dense models, it also notes diminishing benefits as base model sizes grow. \citet{ludziejewski2024scaling} challenge this conclusion, attributing the diminished returns partly to the fixed number of training tokens across models and constant expert sizes. By introducing "granularity" and adjusting training durations, they demonstrate that MoEs can outperform dense models across any compute budget, debunking the notion of diminishing returns for MoEs with adaptive expert configurations. More recently, \citet{jelassi2024mixture} finds that, on downstream tasks, MoEs scale efficiently with the number of experts (i.e., increasing sparsity) on memorization tasks, but their reasoning capabilities saturate and lag behind dense models on tasks requiring complex reasoning when compared based on total number of parameters.

Another approach by \citet{he2024mixture} explores the benefits of training MoEs with larger numbers of smaller experts rather than the conventional setup of fewer, larger experts. They introduce Parameter Efficient Expert Retrieval (PEER), a novel routing mechanism designed to tackle the computational and optimization challenges that arise when handling a high number of experts, thus enabling efficient scaling of MoE models.

Lastly, \citet{yun2024toward} draws attention to the increased inference costs associated with scaling MoEs by adding experts. While additional experts may not substantially affect training costs, they can inflate inference costs, thereby diminishing deployment efficiency. To address this, the study proposes an over-trained budget allocation strategy, optimizing MoE models for both performance and efficiency in deployment.